\subsection{Using both barriers when the upper slope is small}

\begin{lemma}[An envelope from both linear bounds]\label{fprof:smallslope}
Let \(0<a\le1/4\), \(a\le b\le1/2\), and let
\(g:[0,N]\to\mathbb R\) be 1-Lipschitz with \(ax\le g(x)\le bx\).
For \(3N/4\le n_0<N\), an admissible decreasing chain from \(n_0\)
through \(N/2\) to zero can be chosen with
\(O(1+\log(N/(N-n_0)))\) edges and total cost at most
$$
 N\max\{\min(E(a,b),F(a,b)),M(a,b)\},
$$
where
$$
 E(a,b)=\frac{1-2a-2a^2+(2+a)b}{6(1+a)},\qquad
 M(a,b)=\frac{3b}{8}+\frac3{16}-\frac{3a}{4},
$$
and
\begin{equation}\label{fprof:small-F}
 F(a,b)=\frac{b-a}{2(1+a)}+\frac{1-2a}{4}
 -\frac{(1-2b)(1-a)(1+b)}{8(1-b)(1+a)}.
\end{equation}
The mesh and approximate-barrier errors have the same
\(O(KT+Ke_0N)\) bounds as in Lemma~\ref{fprof:reflection}.
\end{lemma}

\begin{proof}
Scale to \(N=1\). Choose a global minimum \(t\) of \(g\) on
\([1/2,1]\), and put \(m=g(t)\), \(v=g(1/2)\), and \(e=g(1)\).
When \(t>3/4\), the late-minimum argument in
Lemma~\ref{fprof:reflection} gives \(M(a,b)\). We therefore assume
\(t\in[1/2,3/4]\).

Set
$$
 z_* =\frac{t-m}{1-b},\qquad
 H(z)=\frac z2+\left(b-\frac12\right)z_*.
$$
Since \(m\le bt\), one has \(z_*\ge t\). On \([t,1]\), the two
available upper bounds give
$$
 g(z)\le\min\{bz,m+z-t\}\le H(z).
$$
For the last inequality, if \(z\le z_*\), then
\(H(z)-(m+z-t)=(z_*-z)/2\ge0\). If \(z\ge z_*\), then
\(H(z)-bz=(1/2-b)(z-z_*)\ge0\). This also covers \(z_*>1\).

For \(x\ge(1+t)/2\), define
$$
 \Lambda(x)=H(2x-1)-g(x).
$$
This is nondecreasing, and almost everywhere
$$
 (-g')_+\le\frac{1-g'}2=\frac{\Lambda'}2.
$$
Use maximal admissible downward jumps in the positive-potential region,
stopping if the endpoint reaches \((1+t)/2\). Their costs sum to at most
\(\Lambda(1)/2\). If the construction first reaches the region
\(\Lambda\le0\) above that stopping threshold, then
$$
 g(2n-1)\le H(2n-1)\le g(n)
$$
at the current endpoint \(n\). Choosing the leftmost minimum on
\([2n-1,n]\) gives an admissible zero-cost edge. Continue until the
endpoint becomes at most \((1+t)/2\). Every such endpoint is at least
\(t\), so the next jump to \(t\) is admissible and costs zero.
If \(n_0\le(1+t)/2\), that direct zero-cost jump suffices from the start.

The construction is first made on a grid. Consecutive zero-cost intervals
can be merged whenever their union is admissible. Every two unmergeable
intervals more than double the complementary depth; the positive-cost
maximal jumps also double it except at their last step. Thus the number
of retained edges has the stated logarithmic bound. Grid crossing errors
are \(O(T)\), and a compactness limit of the bounded-length endpoint
lists proves the continuous assertion. This is the same affine-envelope
construction as before, now stopped at \((1+t)/2\); it does not require
\(\Lambda((1+t)/2)\le0\).

Since \(H(1)\ge e\), the prefix cost is bounded by the nonnegative number
$$
 \frac{H(1)-e}{2}
 =\frac{1/2+(b-1/2)z_*-e}{2}.
$$
The jump from \(t\) to \(1/2\) is admissible and costs \(v-m\).
Combining this construction with the earlier cone construction gives
a cost at most
\begin{equation}\label{fprof:small-scalar}
 v-m+\min\left\{
 \frac{1-t+m-a}{3},\quad
 \frac{1/2-a-k(t-m)}2
 \right\},\qquad
 k=\frac{1/2-b}{1-b}.
\end{equation}
The subsequent edge to zero has zero cost.

It remains to maximize \eqref{fprof:small-scalar} under
$$
 m\ge at,\qquad v\le\min\{b/2,m+t-1/2\}.
$$
For fixed \(t\), replace \(v\) by its upper bound. Both expressions
inside the resulting minimum are nondecreasing in \(m\) below
\(b/2-t+1/2\), and decreasing above it. Hence both attain their maximum
at the same scalar value
$$
 m=\max\{at,b/2-t+1/2\}.
$$
Write \(t_*=(1+b)/(2(1+a))\). For \(t\le t_*\), the two resulting
expressions increase with slopes \(1/3\) and \(1-k\); for
\(t\ge t_*\), they decrease with slopes
\(-(1+2a)/3\) and \(-a-k(1-a)/2\). The point \(t_*\) belongs to
\([1/2,3/4]\), since \(a\le b\le1/2\). Their minimum is therefore
maximized at \(t=t_*\), \(m=at_*\). Substitution gives exactly
\(\min(E(a,b),F(a,b))\).

For approximate barriers, clamp between the exact linear barriers and
interpolate grid values. The uniform change is \(O(e_0N+T)\), and every
edge cost changes by at most twice that quantity. The classification into
early and late cases is applied to the clamped profile itself. This proves
the asserted stability without requiring the minimum location to be
preserved by clamping.
\end{proof}

\subsection{The additional packing interval}

For \(3/16\le a\le1/5\), put \(c=13a+12a^2\), and define
\begin{equation}\label{fprof:small-bF}
 b_F(a)=\frac{3+c-\sqrt{(3+c)^2+8(1+a)(1-c)}}{4(1+a)}.
\end{equation}
Let \(a_L\) be the unique root in \((3/16,19/100)\) of
$$
 16a^3+6a^2-7a+1=0,
$$
and let \(a_J\) be the unique root in \((19/100,1/5)\) of
$$
 32a^4-28a^3-75a^2-5a+4=0.
$$
In particular, \(a_L<a_J\). Their approximate values are
$$
 a_L=0.188783591161\ldots,\qquad
 a_J=0.195322470969\ldots.
$$

\begin{proposition}[Additional small-slope balance]\label{fprof:small-balance}
For \(a_L\le a\le a_J\) and \(a\le b<b_F(a)\),
$$
 \max\{\min(E(a,b),F(a,b)),M(a,b)\}<a.
$$
Moreover,
$$
 b_F(a)>\max\left\{2a,\frac{8a^2+8a-1}{a+2}\right\}
 \qquad(a_L<a<a_J).
$$
At the endpoints,
$$
 b_F(a_L)=2a_L,\qquad
 b_F(a_J)=\frac{8a_J^2+8a_J-1}{a_J+2}.
$$
\end{proposition}

\begin{proof}
Multiplying \(F-a\) by its positive denominator shows
$$
 F(a,b)<a\quad\Longleftrightarrow\quad G(a,b)<0,
$$
where
$$
 G(a,b)=-2(1+a)b^2+(3+13a+12a^2)b+1-13a-12a^2.
$$
For \(0\le b\le1/2\),
$$
 \partial_bG\ge1+11a+12a^2>0.
$$
In the stated range of \(a\),
$$
 G(a,a)=(a-1)(a+1)(10a-1)<0,\qquad
 G(a,1/2)=2-7a-6a^2>0.
$$
Thus \(b_F(a)\), as given by \eqref{fprof:small-bF}, is the unique zero
in \((a,1/2)\), and \(b<b_F(a)\) implies \(F<a\).

For the late bound, \(M<a\) is equivalent to
\(b<R(a):=14a/3-1/2\). On \([3/16,1/5]\), this comparison point is
in \([0,1/2]\), and
$$
 G(a,R(a))=\frac{112a^3+76a^2+30a-9}{9}.
$$
The last expression is increasing for positive \(a\) and equals
\(1/256\) at \(a=3/16\). Therefore \(G(a,R(a))>0\), whence
\(b_F(a)<R(a)\). This proves the strict combined balance.

Next,
$$
 G(a,2a)=16a^3+6a^2-7a+1.
$$
The polynomial on the right is strictly decreasing on
\([3/16,1/5]\); its values at \(3/16\) and \(19/100\) are
\(1/256\) and \(-457/125000\). This proves the existence and uniqueness
of \(a_L\), the first endpoint identity, and \(b_F(a)>2a\) for
\(a_L<a<a_J\).

Finally, with \(b_E(a)=(8a^2+8a-1)/(a+2)\), expansion gives
$$
 (a+2)^2G(a,b_E(a))
 =-(a+1)(32a^4-28a^3-75a^2-5a+4).
$$
The quartic in parentheses is strictly decreasing on \([0,1/4]\), since
its derivative is
$$
 128a^3-84a^2-150a-5
 \le-52a^2-150a-5<0.
$$
Its values at \(19/100\) and \(1/5\) are
\(600471/3125000\) and \(-108/625\), proving the asserted unique
root \(a_J\). Thus \(G(a,b_E(a))<0\) for \(a<a_J\) in the present
parameter interval. Since \(a<b_E(a)<1/2\) there, monotonicity of
\(G\) implies \(b_E(a)<b_F(a)\). Equality at \(a_J\) proves the
second endpoint identity.
\end{proof}

\begin{proposition}[Packing consequence on the additional interval]
\label{fprof:small-packing}
Let \(E\subset\mathbb R^2\) be Borel, and write
\(d=\dim_{\mathrm H}E\), \(D=\dim_{\mathrm P}E\). The finite-profile
transfer in this manuscript, with Lemma~\ref{fprof:smallslope} as its
chain input, gives
$$
 1+a_L<d<1+a_J,\qquad D<1+b_F(d-1)
 \quad\Longrightarrow\quad
 \exists y\in E:\ |\Delta_y(E)|>0.
$$
\end{proposition}

\begin{proof}
Proposition~\ref{fprof:small-balance} gives a strict profile margin at
\(a=d-1\), \(b=D-1\). Since the inequalities are strict and the
expressions are continuous, choose a Frostman exponent \(s<d\) and an
upper covering exponent \(u>D\), sufficiently close to \(d,D\), that
the same margin holds at \(a=s-1\), \(b=u-1\).
Lemma~\ref{fprof:smallslope} supplies bounded-length chains with that
cost and the required regularization errors. Applying the finite-profile
transfer with this chain bound, and the same compact selection and
separation argument, proves the stated positive-length conclusion.
\end{proof}

For example, take \(a=19/100\) and \(b=761/2000\). Then
$$
 F(a,b)=\frac{74596147}{393176000},\qquad
 M(a,b)=\frac{3003}{16000},\qquad
 a-\max(F(a,b),M(a,b))=\frac{107293}{393176000}>0.
$$
Thus the additional interval includes the dimension pair
\(d=1.19\), \(D=1.3805\), while \(2d-1=1.38\).

Combining this with the preceding sufficient bounds gives the continuous
piecewise cutoff
\begin{equation}\label{fprof:smallslopecurve}
 B_{\rm refined}(d)=\begin{cases}
 2d-1,&1<d\le1+a_L,\\[3pt]
 1+b_F(d-1),&1+a_L<d<1+a_J,\\[3pt]
 \dfrac{d(8d-7)}{d+1},&1+a_J\le d\le5/4.
 \end{cases}
\end{equation}
Both joins agree by Proposition~\ref{fprof:small-balance}. As elsewhere,
the distance conclusion requires the separate analytic transfer; the
elementary chain estimates alone are not a proof about distance sets.
